%% Elemento obrigatório (Figura 15).
%% Consiste em uma tradução do resumo em português para uma
%% língua estrangeira (em inglês, ABSTRACT; em espanhol, RESUMEN;
%% em francês, RÉSUMÉ), em um único parágrafo, seguido das palavras-
%% -chave representativas do conteúdo do trabalho, na língua estrangeira
%% escolhida.
%% O resumo em outra língua também é precedido pela referência
%% do trabalho, substituindo-se o título em português pelo título na língua
%% estrangeira adotada.
%% No caso de teses, é possível incluir dois resumos em língua es-
%% trangeira.
%% A apresentação gráfica e a ordem dos elementos seguem a mes-
%% ma orientação do resumo em português.

\begin{resumo}[Abstract]
\begin{otherlanguage*}{english}
\begin{SingleSpace}

\noindent
\begin{flushleft}
\entradaAutor{}. \textit{\englishTitle{}}. \the\year. \pageref{LastPage} p. Undergraduate Thesis (Bachelor's Degree in Computer Engineering) - Instituto Politécnico, Universidade do Estado do Rio de Janeiro, Nova Friburgo, \the\year.
\end{flushleft}
\vspace{\onelineskip}

\setlength{\parindent}{1.3cm}

The failures of tailings dams in Mariana (2015) and Brumadinho (2019)
exposed critical weaknesses in the monitoring and warning systems of the
Brazilian mining area, resulting in hundreds of deaths and
environmental damage of historic proportions. In this context, this
work proposes the development of an automated monitoring and alert
system for mining tailings dams, based on IoT technology, with the
objective of continuously measuring soil moisture, integrating external
meteorological data and calculating a real-time geotechnical risk index.
The system comprises an embedded hardware prototype built on the
ESP8266 NodeMCU microcontroller with a resistive hygrometer sensor,
indicator LEDs and a buzzer, programmed in C++ using the Arduino IDE.
The firmware integrates the BNDMET/DECEA and OpenWeatherMap
meteorological APIs to retrieve historical precipitation data and
rainfall forecasts, transmitting readings to the backend at adaptive
intervals according to the calculated risk level. The integrated risk
assessment model, developed by the author, calculates a risk factor
through the weighted sum of seven variables expressing geotechnical
conditions --- soil moisture and rate of change --- and meteorological
conditions --- accumulated precipitation over multiple time windows,
rainfall forecast and atmospheric pressure ---, with the rainfall
thresholds calibrated based on \citeonline{rico2008mine} and on the
criteria of the Alerta Rio System \cite{sistemaalertario2025}, and an
amplification mechanism triggered when high soil saturation coincides
with forecast rainfall. The software platform adopts a three-layer client-server
architecture: embedded firmware, a backend developed in Node.js with
TypeScript and a PostgreSQL database with the TimescaleDB extension for
time-series data, and a React/Next.js frontend providing a real-time
monitoring interface, reading history, user management and an alert
center. The system classifies risk into three levels (green, yellow, red), 
signalled physically by the LEDs and buzzer on the embedded
device and displayed on the web platform. The administrator has access to
an alert center that enables the manual dispatch of bulk e-mail
notifications to registered users, accompanied by the corresponding
analysis cycle report. The proposed solution demonstrates that it is
possible to develop, using low-cost hardware and open-source technologies,
a functional system for continuous geotechnical risk monitoring,
which may include the measurement and use of other variables, 
such as the water table level and soil pore pressure (piezometer) and the velocity, 
direction, and depth of landslides (inclinometer),
contributing to accident prevention in mining tailings dams.

\vspace{\onelineskip}
\noindent Keywords: Alert systems; ESP8266; Geotechnical monitoring; Internet of Things; Tailings dams.

\end{SingleSpace}
\end{otherlanguage*}
\end{resumo}